\documentclass[11pt]{article}
\usepackage[a4paper,margin=1in]{geometry}
\usepackage{fontspec}
\usepackage{amsmath,amssymb,amsthm}
\usepackage{booktabs,longtable,array}
\usepackage{enumitem}
\usepackage{xcolor}
\usepackage[colorlinks=true,linkcolor=blue,citecolor=blue,urlcolor=blue]{hyperref}
\setlist{nosep}
\setlength{\parskip}{0.6em}
\setlength{\parindent}{0pt}

\title{Topos-Theoretic Frameworks for Neural Networks and LLMs}
\author{}
\date{\today}

\begin{document}
\maketitle


\begin{abstract}
This article surveys a recent line of work whose common claim is not that topos theory has already produced a mature engineering paradigm for artificial intelligence, but rather that certain neural-network and large-language-model structures can be represented, completed, or semantically organized inside categorical and topos-theoretic frameworks. We give a unified preliminary account of the categorical language used in these works, beginning with categories, functors, natural transformations, limits, adjunctions, presheaves, Yoneda, sheaves, sites, Grothendieck topologies, Grothendieck topoi, internal logic, and stacks. We then explain how this language is used in several representative proposals: Lafforgue's programmatic use of Grothendieck topoi for AI, Belfiore and Bennequin's topos-and-stack framework for deep neural networks, the withdrawn but suggestive proposal on the topos of transformer networks, Mahadevan's topos-theoretic proposal for generative AI and LLMs, and the 2025 survey by Jia, Peng, Yang, and Chen. The emphasis is on what is mathematically constructed or claimed, what is merely programmatic, and which notions from topos theory are actually doing work.


\end{abstract}

\section{Introduction}

Modern neural networks are usually described as parametrized maps


\[
N_\theta:X\longrightarrow Y
\]

or as composites of layers


\[
N_\theta
=f_L\circ f_{L-1}\circ\cdots\circ f_1.
\]

This description is effective for optimization, implementation, and empirical evaluation, but it says relatively little about several structural questions:


\begin{enumerate}
\item How should local representations, layerwise data, modules, symmetries, and
invariances be organized into global semantic structures?

\item Can a neural network be represented as an object in a category rich enough to
support local-to-global reasoning?

\item Can the expressive power of transformers or LLMs be described through
categorical completions, function objects, subobjects, or internal logic?

\item Can model composition be governed by universal constructions such as
pullbacks, pushouts, equalizers, coequalizers, and exponentials?

\end{enumerate}

Topos theory is relevant to these questions because a Grothendieck topos is, roughly, a universe of sheaves on a site:


\[
\mathcal{E}\simeq \mathbf{Sh}(\mathcal{C},J).
\]

Here \(\mathcal{C}\) is a category of contexts, \(J\) is a Grothendieck topology encoding which families count as covers, and \(\mathbf{Sh}(\mathcal{C},J)\) is the category of sheaves satisfying a local-to-global gluing condition. Such a category behaves simultaneously like a generalized space, a generalized universe of sets, and a universe with its own internal logic.


The works discussed below should be read with care. Some contain precise categorical constructions; some are programmatic; some are speculative; and one important transformer-related proposal is marked withdrawn on arXiv. Their common contribution is not a uniform theorem about test accuracy or generalization, but a family of mathematical attempts to place neural architectures in richer semantic environments.


\section{Unified Notation}

We use the following notation throughout.


\begin{longtable}{>{\raggedright\arraybackslash}p{0.24\textwidth}>{\raggedright\arraybackslash}p{0.66\textwidth}}
\toprule
Symbol & Meaning \\
\midrule
\endhead
\(\mathcal{C},\mathcal{D}\) & categories, often small when presheaves are formed \\
\(A,B,C,U,V\) & objects of a category \\
\(f:A\to B\) & a morphism \\
\(\mathbf{Set}\) & category of sets and functions \\
\(\mathbf{PSh}(\mathcal{C})\) & presheaf category \([\mathcal{C}^{op},\mathbf{Set}]\) \\
\(F,G\) & functors or presheaves \\
\(\eta:F\Rightarrow G\) & natural transformation \\
\(y:\mathcal{C}\to \widehat{\mathcal{C}}\) & Yoneda embedding \\
\(yA\) & representable presheaf \(\mathrm{Hom}_{\mathcal{C}}(-,A)\) \\
\(J\) & Grothendieck topology on \(\mathcal{C}\) \\
\((\mathcal{C},J)\) & a site \\
\(\mathbf{Sh}(\mathcal{C},J)\) & category of sheaves on a site \\
\(\mathcal{E},\mathcal{F}\) & topoi \\
\(\Omega\) & subobject classifier \\
\(N_\theta\) & neural network with parameters \(\theta\) \\
\(h_\ell\) & hidden representation at layer \(\ell\) \\
\(T_\theta\) & transformer or LLM-type model \\
\bottomrule
\end{longtable}

For neural networks we use a deliberately light notation. A feedforward network is a composite


\[
h_{\ell+1}=f_{\ell,\theta_\ell}(h_\ell),
\qquad
N_\theta=f_{L,\theta_L}\circ\cdots\circ f_{1,\theta_1}.
\]

A transformer block may be written schematically as


\[
H_{\ell+1}
=\mathrm{FFN}_\ell\bigl(\mathrm{Attn}_\ell(H_\ell)\bigr),
\]

where \(H_\ell\) is a sequence of hidden states. The exact algorithmic details are not central here; the categorical question is how such maps, layers, heads, states, symmetries, and tasks may be organized as categorical objects and morphisms.


\section{Categorical Preliminaries}

\subsection{Categories, Functors, and Natural Transformations}

A category \(\mathcal{C}\) consists of objects, morphisms, identity morphisms, and associative composition. For any objects \(A,B\in\mathcal{C}\), the set of morphisms from \(A\) to \(B\) is denoted


\[
\mathrm{Hom}_{\mathcal{C}}(A,B).
\]

A functor


\[
F:\mathcal{C}\to\mathcal{D}
\]

sends objects to objects and morphisms to morphisms while preserving identities and composition:


\[
F(\mathrm{id}_A)=\mathrm{id}_{F(A)},\qquad
F(g\circ f)=F(g)\circ F(f).
\]

A natural transformation \(\eta:F\Rightarrow G\) between two functors \(F,G:\mathcal{C}\to\mathcal{D}\) assigns to every \(A\in\mathcal{C}\) a morphism


\[
\eta_A:F(A)\to G(A)
\]

such that for every \(f:A\to B\),


\[
G(f)\circ \eta_A=\eta_B\circ F(f).
\]

In the topos-AI literature, functors are used to express structured translation: from geometric data to semantic data, from network components to categorical objects, or from syntactic descriptions to semantic models. Natural transformations express compatibility between such translations.


\subsection{Limits and Universal Properties}

Limits are categorical constructions defined by universal properties. The most important finite limits are terminal objects, products, equalizers, and pullbacks.


A pullback of


\[
A\xrightarrow{f}C\xleftarrow{g}B
\]

is an object \(A\times_C B\) equipped with projections


\[
p_1:A\times_C B\to A,\qquad p_2:A\times_C B\to B
\]

such that


\[
f\circ p_1=g\circ p_2,
\]

and universal with this property. In \(\mathbf{Set}\),


\[
A\times_C B
=\{(a,b)\in A\times B\mid f(a)=g(b)\}.
\]

In AI language, pullbacks model alignment over a shared semantic target. If \(I\to S\) and \(T\to S\) map image and text representations into a shared semantic space \(S\), then \(I\times_S T\) models semantically compatible image-text pairs. This is not yet topos theory, but it is the kind of universal construction that later appears inside topoi.


\subsection{Adjunctions}

An adjunction


\[
L\dashv R
\]

between \(L:\mathcal{C}\to\mathcal{D}\) and \(R:\mathcal{D}\to\mathcal{C}\) is a natural bijection


\[
\mathrm{Hom}_{\mathcal{D}}(LC,D)
\cong
\mathrm{Hom}_{\mathcal{C}}(C,RD).
\]

Left adjoints often perform free constructions or completions; right adjoints often forget or restrict structure. Examples include:


\[
\text{free group }F:\mathbf{Set}\to\mathbf{Grp}
\quad\dashv\quad
U:\mathbf{Grp}\to\mathbf{Set}.
\]

Adjunctions are central in topos theory. Sheafification is a left adjoint to the inclusion of sheaves into presheaves. Geometric morphisms are adjoint pairs of functors. Existential and universal quantification in internal logic can also be expressed via adjoints.


\subsection{Presheaves and Yoneda}

For a category \(\mathcal{C}\), a presheaf is a functor


\[
F:\mathcal{C}^{op}\to\mathbf{Set}.
\]

The category of presheaves is


\[
\widehat{\mathcal{C}}
=\mathbf{PSh}(\mathcal{C})
=[\mathcal{C}^{op},\mathbf{Set}].
\]

For \(A\in\mathcal{C}\), the representable presheaf \(yA\) is


\[
yA=\mathrm{Hom}_{\mathcal{C}}(-,A).
\]

The Yoneda lemma states that for every presheaf \(F\),


\[
\mathrm{Nat}(yA,F)\cong F(A).
\]

Thus an element of \(F(A)\) is equivalent to a natural way of producing \(F\)-data from every generalized element \(X\to A\). The Yoneda embedding


\[
y:\mathcal{C}\to\widehat{\mathcal{C}}
\]

is fully faithful:


\[
\mathrm{Hom}_{\mathcal{C}}(A,B)
\cong
\mathrm{Nat}(yA,yB).
\]

This is why presheaf categories are natural completions of a starting category: they contain the original category fully faithfully and add generalized context-dependent objects.


\subsection{Sites, Sheaves, and Sheafification}

A Grothendieck topology \(J\) on a category \(\mathcal{C}\) specifies, for each object \(U\), which families


\[
\{U_i\to U\}_{i\in I}
\]

count as covers. The pair \((\mathcal{C},J)\) is called a site.


A presheaf \(F:\mathcal{C}^{op}\to\mathbf{Set}\) is a sheaf if compatible local sections over every covering family glue uniquely to a global section. In the classical topological case, this says: if


\[
U=\bigcup_i U_i
\]

and \(s_i\in F(U_i)\) satisfy


\[
s_i|_{U_i\cap U_j}=s_j|_{U_i\cap U_j},
\]

then there exists a unique \(s\in F(U)\) such that


\[
s|_{U_i}=s_i.
\]

The sheafification of a presheaf \(F\) is a sheaf \(aF\) equipped with a map \(F\to iaF\), where \(i:\mathbf{Sh}(\mathcal{C},J)\hookrightarrow \mathbf{PSh}(\mathcal{C})\) is the inclusion. It is characterized by the adjunction


\[
a\dashv i.
\]

Concretely, sheafification repairs two failures:


\begin{enumerate}
\item local equality should imply global equality;
\item compatible local data should have a global gluing.
\end{enumerate}

This is one of the main reasons sheaves are relevant to neural systems: they formalize the passage from local descriptions to coherent global descriptions.


\subsection{Grothendieck Topoi}

A Grothendieck topos is a category equivalent to the category of sheaves on a site:


\[
\mathcal{E}\simeq\mathbf{Sh}(\mathcal{C},J).
\]

It has three complementary interpretations:


\begin{enumerate}
\item a generalized space;
\item a generalized universe of sets varying over contexts;
\item a universe with an internal logic.
\end{enumerate}

Examples include:


\[
\mathbf{Set}\simeq\mathbf{Sh}(*),
\]

the sheaf topos \(\mathbf{Sh}(X)\) of a topological space \(X\), and every presheaf category


\[
[\mathcal{C}^{op},\mathbf{Set}].
\]

In the AI works discussed below, a Grothendieck topos is often used as a semantic environment: a place where local structures, invariances, data, and logical assertions can be organized coherently.


\subsection{Subobject Classifiers and Internal Logic}

An elementary topos has finite limits, exponentials, and a subobject classifier \(\Omega\). The subobject classifier generalizes the two-element set of truth values. In \(\mathbf{Set}\), a subset \(S\subseteq X\) is classified by its characteristic function


\[
\chi_S:X\to\{0,1\}.
\]

In a topos \(\mathcal{E}\), a subobject \(S\hookrightarrow X\) is classified by a morphism


\[
\chi_S:X\to\Omega.
\]

Thus predicates become subobjects, and truth values live in \(\Omega\). A topos therefore supports an internal logic. This matters for AI proposals that seek to represent semantic assertions, concept membership, interpretability claims, or contextual truth inside a categorical universe rather than as external Boolean labels.


\subsection{Stacks}

A sheaf glues set-valued local data. A stack glues category-valued local data, including not only objects but also morphisms or isomorphisms between objects. Very roughly:


\[
\text{sheaf}:\text{ local objects glue},
\qquad
\text{stack}:\text{ local objects and equivalences glue}.
\]

Stacks are invoked in the DNN literature when local representations are not expected to match strictly but only up to isomorphism, symmetry, gauge transformation, or change of coordinates. This is particularly natural for neural-network layers, heads, and feature spaces, where equivalent representations need not be literally equal.


\section{Minimal AI Background}

For the purposes of this article, a feedforward DNN is a parametrized composite


\[
N_\theta=f_{L,\theta_L}\circ\cdots\circ f_{1,\theta_1}.
\]

A CNN adds locality and translation-equivariant filters; an RNN or LSTM adds stateful sequence processing; a transformer uses attention mechanisms to compute contextualized token representations. A large language model is a large-scale transformer trained to predict or model token distributions.


The topos-theoretic works surveyed below do not usually depend on the numerical details of training. Instead, they abstract neural architectures as structured systems of components, maps, states, parameters, symmetries, and semantic relations.


\section{Lafforgue: Grothendieck Topoi as Semantic Completions for AI}

Laurent Lafforgue's 2022 lecture note, "Some Possible Roles for AI of Grothendieck Topos Theory," is programmatic. Its central motivation is that neural networks manipulate numerical functions, but mathematical meaning often comes from structured geometric or logical origins. The proposed role of Grothendieck topoi is to preserve and enrich such meaningful origins.


In categorical terms, one begins with a category \(\mathcal{C}\) of meaningful objects, such as geometric shapes, symbolic structures, or knowledge objects. Through the Yoneda embedding,


\[
y:\mathcal{C}\to\widehat{\mathcal{C}},
\]

the original category is embedded into a presheaf universe. By imposing a Grothendieck topology \(J\) and passing to sheaves,


\[
\mathbf{Sh}(\mathcal{C},J),
\]

one obtains a richer category in which local-to-global structure, generalized functions, and semantic invariants can be handled.


The conclusion is not an implemented learning algorithm. It is a research program:


\[
\text{meaningful category}
\longrightarrow
\text{presheaf completion}
\longrightarrow
\text{sheaf/topos completion}.
\]

The relevant preliminary notions are Yoneda, sheafification, sites, Grothendieck topologies, and Grothendieck topoi. Lafforgue's proposal uses topos theory as a semantic completion mechanism, not as a direct replacement for gradient-based learning.


\section{Belfiore and Bennequin: Topoi and Stacks of Deep Neural Networks}

Jean-Claude Belfiore and Daniel Bennequin's paper "Topos and Stacks of Deep Neural Networks" is one of the most ambitious direct proposals connecting DNNs to Grothendieck topoi and stacks.


Its central claim can be summarized as:


\[
\boxed{\parbox{0.86\textwidth}{\centering
Known deep neural networks can be represented as objects in canonical
Grothendieck topoi, and their invariance structures can be described by stacks.}}
\]

The idea is that a DNN is not merely a function


\[
N_\theta:X\to Y,
\]

but a structured object involving layers, states, parameters, transformations, symmetries, and learning dynamics. These structures can be organized into categories and sites. The associated Grothendieck topoi provide global semantic environments for the network, while stacks encode invariance and equivalence data.


Several preliminary notions are directly relevant:


\begin{itemize}
\item \textbf{Presheaves and sheaves.} Layerwise or local network data may be viewed as
data assigned to contexts, with restriction and gluing.

\item \textbf{Grothendieck topoi.} The network is placed in a sheaf-theoretic universe.
\item \textbf{Internal logic.} Semantic assertions about network behavior can be
interpreted inside the topos.

\item \textbf{Stacks.} CNNs, LSTMs, or other structured architectures may involve
invariances and equivalences that should glue only up to isomorphism.

\end{itemize}

This framework should be read as a structural encoding proposal. It does not by itself prove that a DNN generalizes, nor does it give a standard training algorithm. Rather, it proposes a categorical semantics in which DNNs and their symmetries may be represented more faithfully than by a bare composite of functions.


\section{The Topos of Transformer Networks: A Withdrawn but Suggestive Proposal}

Villani and McBurney's "The Topos of Transformer Networks" attempts to relate transformer expressivity to topos-theoretic completion. The arXiv page marks the paper as withdrawn and requiring major revision, so its claims must be treated as unstable.


The proposed picture is roughly as follows. Many conventional neural-network architectures are said to embed in a pretopos-like category of piecewise-linear functions. Transformers, by contrast, are argued to require a richer topos completion, reflecting their contextual, higher-order, and attention-based structure.


The relevant preliminary notions are:


\begin{itemize}
\item \textbf{Pretopos.} A category with finite limits, disjoint finite coproducts, and
effective quotients in a suitable sense.

\item \textbf{Topos completion.} A process of embedding a weaker categorical environment
into a richer topos-like one.

\item \textbf{Exponentials and internal logic.} Transformers are interpreted as needing
richer function-object and higher-order expressive capacity.

\end{itemize}

Because the paper is withdrawn, the responsible conclusion is not that transformers have been rigorously characterized by a stable topos theorem. Rather, this work points toward a research question:


\[
\parbox{0.86\textwidth}{\centering
Does attention-based contextual computation require categorical
structure beyond ordinary first-order neural-network semantics?}
\]

This is a useful idea, but not currently a settled result.


\section{Mahadevan: Topos Theory for Generative AI and LLMs}

Mahadevan's "Topos Theory for Generative AI and LLMs" proposes that categories of LLM-like systems may possess topos-theoretic structure and that universal constructions can guide the design of compositional AI architectures.


The claimed direction is:


\begin{enumerate}
\item treat LLMs or generative models as objects or morphisms in a category;
\item show the resulting category has limits, colimits, exponentials, and
subobject-classifier-like structure;

\item use universal constructions to design or analyze model composition.
\end{enumerate}

For example:


\begin{itemize}
\item pullbacks model semantic alignment;
\item pushouts model modular fusion;
\item equalizers model agreement between model outputs;
\item coequalizers model quotienting by behavioral equivalence;
\item exponentials model conditional or higher-order model spaces;
\item subobject classifiers model predicates or classifiers over model behavior.
\end{itemize}

This work therefore uses almost every major preliminary notion above: limits, colimits, exponentials, subobjects, truth values, and categorical universes. The main caution is that the construction depends heavily on how the category of models is defined. Without a precise category and precise morphisms, the statement "LLMs form a topos" is a guiding slogan rather than a mathematical theorem.


\section{Jia-Peng-Yang-Chen Survey}

The survey "Category-Theoretical and Topos-Theoretical Frameworks in Machine Learning" classifies categorical machine learning into several broad tracks:


\begin{enumerate}
\item gradient-based learning;
\item probability-based learning;
\item invariance and equivalence-based learning;
\item topos-based learning.
\end{enumerate}

The first three tracks are closer to established applied category theory: Cartesian differential categories, lenses, Markov categories, categorical probability, and equivariant learning. The fourth track collects more exploratory work involving Grothendieck topoi, stacks, sheaves, and internal logic.


The survey is useful as a map, not as a single theorem. Its main contribution is to identify where topos-theoretic proposals sit relative to more mature categorical ML frameworks.


\section{Comparison of the Works}

\begin{longtable}{>{\raggedright\arraybackslash}p{0.18\textwidth}>{\raggedright\arraybackslash}p{0.34\textwidth}>{\raggedright\arraybackslash}p{0.25\textwidth}>{\raggedright\arraybackslash}p{0.13\textwidth}}
\toprule
Work & Main construction or claim & Topos notions used & Status \\
\midrule
\endhead
Lafforgue & Grothendieck topoi as semantic completions for meaningful AI objects & Yoneda, sheaves, sites, topoi & programmatic \\
Belfiore-Bennequin & DNNs as objects in canonical Grothendieck topoi; invariance via stacks & Grothendieck topoi, stacks, internal logic & ambitious structural framework \\
Transformer topos paper & transformers require topos completion beyond pretopos-like neural categories & pretopos, topos completion, exponentials & withdrawn; research pointer \\
Mahadevan & LLM categories may form topoi; universal constructions guide AI composition & limits, colimits, exponentials, subobject classifier & speculative proposal \\
Jia-Peng-Yang-Chen & survey of categorical and topos-theoretic ML & broad survey & literature map \\
\bottomrule
\end{longtable}

The common theme is not an empirical theorem but a structural thesis:


\[
\parbox{0.86\textwidth}{\centering
Neural systems should be studied as objects in categories rich enough to
encode context, locality, semantics, invariance, and internal logic.}
\]

\section{What Is Actually Proved?}

It is useful to distinguish four levels of claim.


\subsection{Level 1: Established Topos Theory}

The definitions and theorems about presheaves, sheaves, sheafification, Grothendieck topoi, internal logic, and geometric morphisms are classical mathematics. These are reliable.


\subsection{Level 2: Structural Representation Claims}

Some works claim that certain neural architectures can be represented in topos-theoretic or stack-theoretic settings. These are mathematical modeling claims. Their value depends on whether the representation preserves important network structure rather than merely encoding everything trivially.


\subsection{Level 3: Semantic Interpretation Claims}

Claims about meaning, interpretability, contextual truth, and internal logic are philosophically and mathematically suggestive, but require careful definitions of the relevant predicates, contexts, covers, and semantic objects.


\subsection{Level 4: Algorithmic or Empirical Claims}

Very few of these works establish direct performance improvements, generalization bounds, or practical algorithms comparable to standard ML papers. Such claims remain largely open.


\section{Research Outlook}

The most plausible near-term research directions are:


\begin{enumerate}
\item \textbf{Sheaf-theoretic consistency for model outputs.} Use local-to-global
structures to measure whether local explanations or predictions glue.

\item \textbf{Categorical semantics for modular model composition.} Use pullbacks,
pushouts, equalizers, and coequalizers to define principled composition.

\item \textbf{Topos-inspired knowledge representation.} Treat contextual truth and
evidence as internal logical structure rather than global Boolean labels.

\item \textbf{Stack-like invariance in representations.} Use stacks or groupoids to
describe equivalences between learned representations.

\item \textbf{Precise categories of LLMs.} Define objects and morphisms carefully enough
that claims about limits, exponentials, and subobject classifiers become genuine theorems.

\end{enumerate}

\section{Conclusion}

Topos-theoretic approaches to neural networks and LLMs are best understood as a developing semantic and structural program. Their central claim is that neural-network architectures are not merely numerical functions but structured objects involving local data, context dependence, invariance, and logical content. Grothendieck topoi, sheaves, stacks, and internal logic provide a language for organizing these features.


The responsible reading is therefore twofold. The categorical preliminary material is stable and powerful; its application to AI is promising but uneven. The most mature neighboring area is sheaf and topological deep learning. Direct topos-based accounts of DNNs, transformers, and LLMs remain exploratory, but they identify a coherent mathematical ambition: to move from black-box parametrized maps to structured semantic universes in which local computation, global meaning, and contextual truth can be studied together.


\section{References}

\begin{itemize}
\item Jean-Claude Belfiore and Daniel Bennequin,
\href{https://arxiv.org/abs/2106.14587}{Topos and Stacks of Deep Neural Networks}.

\item Laurent Lafforgue,
\href{https://www.laurentlafforgue.org/Expose_Lafforgue_topos_AI_ETH_sept_2022.pdf}{Some Possible Roles for AI of Grothendieck Topos Theory}.

\item Mattia Jacopo Villani and Peter McBurney,
\href{https://arxiv.org/abs/2403.18415}{The Topos of Transformer Networks}. The arXiv record marks this paper as withdrawn and requiring major revision.

\item Sridhar Mahadevan,
\href{https://arxiv.org/abs/2508.08293}{Topos Theory for Generative AI and LLMs}.

\item Yiyang Jia, Guohong Peng, Zheng Yang, and Tianhao Chen,
\href{https://arxiv.org/abs/2408.14014}{Category-Theoretical and Topos-Theoretical Frameworks in Machine Learning: A Survey}; published version: \href{https://doi.org/10.3390/axioms14030204}{Axioms 2025, 14, 204}.

\item Saunders Mac Lane and Ieke Moerdijk, \emph{Sheaves in Geometry and Logic}.
\item Peter Johnstone, \emph{Sketches of an Elephant: A Topos Theory Compendium}.
\item Steve Awodey, \emph{Category Theory}.
\end{itemize}


\end{document}
